% Generated human-review packet template.  Source records, Lean statements,
% and raw-source-to-Spec screening fields are substituted by
% scripts/review_dashboard_packet.py.
% Do not hand-edit generated packets; annotate the PDF or reconcile notes into
% the paper-local human-review log.
\documentclass[10pt]{article}
\usepackage[margin=0.43in]{geometry}
\usepackage{fontspec}
\setmainfont{Latin Modern Roman}
\setmonofont{DejaVu Sans Mono}
\usepackage{hyperref}
\usepackage{amssymb}
\usepackage{fvextra}
\usepackage{enumitem}
\setlength{\parindent}{0pt}
\setlength{\parskip}{0.15em}
\linespread{0.92}
\hypersetup{colorlinks=true,linkcolor=blue,urlcolor=blue}
\DefineVerbatimEnvironment{ReviewVerbatim}{Verbatim}{breaklines=true,breakanywhere=true,fontsize=\scriptsize}
\newcommand{\reviewerbox}{%
  \par\noindent\fbox{\parbox{0.96\linewidth}{\rule{0pt}{0.38in}}}\par
}
\newcommand{\reviewmatch}[1]{%
  \par\noindent\textbf{Reviewer check:} $\square$\; #1\par
  \noindent\emph{If left unchecked, explain the discrepancy or uncertainty below.}\par
}

\begin{document}
\fontsize{9}{10}\selectfont

\begin{center}
{\LARGE Human review packet: MBJG25ProducerFairness}\\[0.35em]
\small Generated from byte-pinned source excerpts, the expanded PaperInterface
specifications, and hash-bound raw-source-to-Spec screening records.
\end{center}

\noindent\textbf{Paper:} MBJG25ProducerFairness\\
\textbf{Paper id:} \texttt{MBJG25ProducerFairness}\\
\textbf{Source version:} not recorded\\
\textbf{Source record:} \texttt{audit/paper\_statement\_map.json}\\
\textbf{Rows:} 28\\
\textbf{Generated:} 2026-08-18\\
\textbf{Source URL:} \href{https://doi.org/10.1609/icwsm.v19i1.35854}{official source}

\noindent\fbox{\parbox{0.96\linewidth}{\textbf{Review status:} 0/28 source-claim screens; 0/5 library prerequisite screens. This is a review worksheet while those direct checks remain pending; it is not a final validation receipt.}}\par



\section*{How to use this packet}
For each source claim, compare the exact source excerpts with the expanded
PaperInterface specification. Mark up this PDF or fill the blank reviewer box in the TeX source;
return the annotated file for reconciliation into the paper-local human-review
log.

\section*{Open the interactive dashboard}
The interactive dashboard is an optional alternative to using this PDF; it
presents the same review surface rather than an additional review requirement.
After forking and cloning (or pulling) EconCSLib, run the following from the
repository root. The cache command is needed only the first time in a new clone.
\begin{ReviewVerbatim}
cd <your-EconCSLib-checkout>
lake exe cache get
python3 scripts/review_dashboard.py --paper MBJG25ProducerFairness --serve
\end{ReviewVerbatim}
Open \url{http://127.0.0.1:8765/} in a browser and keep that terminal running
while reviewing. The dashboard records saved annotations in this paper's
reviewer-owned review record.

\section*{Contents}
\begin{itemize}[leftmargin=1.2em]
\item \hyperlink{library-section-ad8f8e32cbdfebc6}{Semantic library prerequisites (5; not paper claims)}
\begin{itemize}[leftmargin=1.2em]
\item \hyperlink{library-prerequisite-035345cdc23fb9d0}{EconCSLib.Statistics.GlobalMaxAt}
\item \hyperlink{library-prerequisite-4257d258b40d00b7}{EconCSLib.Statistics.GlobalMinAt}
\item \hyperlink{library-prerequisite-1661d59161ee9155}{EconCSLib.Statistics.JensenConcave}
\item \hyperlink{library-prerequisite-aa3d377dcc497061}{EconCSLib.Statistics.JensenConvex}
\item \hyperlink{library-prerequisite-e5c7eedc40c11b7a}{EconCSLib.pmfExp}
\end{itemize}
\item \hyperlink{source-section-f10b5f68e4acf3dc}{Source claims (28)}
\begin{itemize}[leftmargin=1.2em]
\item \hyperlink{source-claim-7a45c0b2fffc539f}{paper\_posterior\_ratingSpec}
\item \hyperlink{source-claim-ec7384519273acb4}{paper\_posterior\_meanSpec}
\item \hyperlink{source-claim-20150c21c10be6e2}{paper\_biasSpec}
\item \hyperlink{source-claim-23bc13a84ad0b29a}{paper\_varianceSpec}
\item \hyperlink{source-claim-6e52f3be6106ad49}{paper\_squared\_biasSpec}
\item \hyperlink{source-claim-9fd08233a5aeb857}{paper\_fixed\_market\_mseSpec}
\item \hyperlink{source-claim-da6d1b14c0df1e3e}{paper\_facing\_fixed\_mse\_decompositionSpec}
\item \hyperlink{source-claim-000171f0bf87b98c}{paper\_facing\_theorem3\_1\_variance\_weak\_decreaseSpec}
\item \hyperlink{source-claim-976fc5a9912d46b4}{paper\_facing\_theorem3\_1\_variance\_strict\_decrease\_interiorSpec}
\item \hyperlink{source-claim-5564d9d725c13f20}{paper\_facing\_theorem3\_1\_squared\_bias\_nondecreasingSpec}
\item \hyperlink{source-claim-2dcf6dbcf162b5d0}{paper\_facing\_theorem3\_2\_squared\_bias\_convex\_in\_qualitySpec}
\item \hyperlink{source-claim-444ef89e21ba4476}{paper\_facing\_theorem3\_2\_squared\_bias\_global\_min\_at\_prior\_meanSpec}
\item \hyperlink{source-claim-3afe5de598f0ed9d}{paper\_facing\_theorem3\_2\_variance\_concave\_in\_qualitySpec}
\item \hyperlink{source-claim-a4656e58047f308c}{paper\_facing\_theorem3\_2\_variance\_global\_max\_at\_halfSpec}
\item \hyperlink{source-claim-6ce0e529b89e57b7}{paper\_facing\_theorem3\_1\_variance\_strict\_decrease\_counterexample\_quality\_zeroSpec}
\item \hyperlink{source-claim-464dc3b105b500a1}{paper\_facing\_theorem3\_1\_variance\_strict\_decrease\_counterexample\_quality\_oneSpec}
\item \hyperlink{source-claim-5e3bedb3cf1d62b0}{paper\_facing\_selection\_rateSpec}
\item \hyperlink{source-claim-430cb3de083df7cf}{paper\_facing\_individual\_producer\_unfairnessSpec}
\item \hyperlink{source-claim-aace03b9ed863bb2}{paper\_facing\_marketplace\_producer\_unfairnessSpec}
\item \hyperlink{source-claim-f5087ebba0ec2e7e}{paper\_facing\_thompson\_sampling\_mechanismSpec}
\item \hyperlink{source-claim-9acd03585d100ede}{paper\_facing\_expected\_regretSpec}
\item \hyperlink{source-claim-7a7d1cda2bb14e73}{paper\_facing\_responsive\_mse\_decompositionSpec}
\item \hyperlink{source-claim-8281b9d563b2f959}{paper\_facing\_eq20\_bayesian\_rating\_equivalenceSpec}
\item \hyperlink{source-claim-48bd38a1dccb3b5f}{paper\_ordinal\_posterior\_ratingSpec}
\item \hyperlink{source-claim-c141e929e91e7be9}{paper\_facing\_eq21\_zero\_prior\_comparisonSpec}
\item \hyperlink{source-claim-c0449d6c0ae69374}{paper\_facing\_k\_sampling\_mechanismSpec}
\item \hyperlink{source-claim-03b433e77eab6bcb}{paper\_facing\_k\_sampling\_one\_is\_argmaxSpec}
\item \hyperlink{source-claim-48e4400514808bc7}{paper\_facing\_k\_sampling\_market\_size\_is\_uniformSpec}
\end{itemize}
\end{itemize}

\clearpage
\hypertarget{library-section-ad8f8e32cbdfebc6}{}
\section*{Material library prerequisites}
\hypertarget{library-prerequisite-035345cdc23fb9d0}{}
\subsection*{\small\ttfamily\raggedright EconCSLib.\allowbreak{}Statistics.\allowbreak{}GlobalMaxAt}
\textbf{Source connection:} no explicit byte-pinned paper source connection is registered.
\paragraph{Lean-expanded library semantic target}
\begin{ReviewVerbatim}
∀ (f : ℝ → ℝ) (x₀ x : ℝ), f x ≤ f x₀
\end{ReviewVerbatim}

\paragraph{Exact Lean library declaration}
\begin{ReviewVerbatim}
def GlobalMaxAt (f : ℝ → ℝ) (x₀ : ℝ) : Prop :=
  ∀ x, f x ≤ f x₀
\end{ReviewVerbatim}

\paragraph{Recorded source-to-library screening}
\textbf{Verdict:} \texttt{not recorded}
\reviewmatch{Matches source input}
\noindent\textbf{Reviewer annotation}\par
\reviewerbox
\clearpage
\hypertarget{library-prerequisite-4257d258b40d00b7}{}
\subsection*{\small\ttfamily\raggedright EconCSLib.\allowbreak{}Statistics.\allowbreak{}GlobalMinAt}
\textbf{Source connection:} no explicit byte-pinned paper source connection is registered.
\paragraph{Lean-expanded library semantic target}
\begin{ReviewVerbatim}
∀ (f : ℝ → ℝ) (x₀ x : ℝ), f x₀ ≤ f x
\end{ReviewVerbatim}

\paragraph{Exact Lean library declaration}
\begin{ReviewVerbatim}
def GlobalMinAt (f : ℝ → ℝ) (x₀ : ℝ) : Prop :=
  ∀ x, f x₀ ≤ f x
\end{ReviewVerbatim}

\paragraph{Recorded source-to-library screening}
\textbf{Verdict:} \texttt{not recorded}
\reviewmatch{Matches source input}
\noindent\textbf{Reviewer annotation}\par
\reviewerbox
\clearpage
\hypertarget{library-prerequisite-1661d59161ee9155}{}
\subsection*{\small\ttfamily\raggedright EconCSLib.\allowbreak{}Statistics.\allowbreak{}JensenConcave}
\textbf{Source connection:} no explicit byte-pinned paper source connection is registered.
\paragraph{Lean-expanded library semantic target}
\begin{ReviewVerbatim}
∀ (f : ℝ → ℝ) (x y lam : ℝ), 0 ≤ lam → lam ≤ 1 → lam * f x + (1 - lam) * f y ≤ f (lam * x + (1 - lam) * y)
\end{ReviewVerbatim}

\paragraph{Exact Lean library declaration}
\begin{ReviewVerbatim}
def JensenConcave (f : ℝ → ℝ) : Prop :=
  ∀ x y lam, 0 ≤ lam → lam ≤ 1 →
    lam * f x + (1 - lam) * f y ≤ f (lam * x + (1 - lam) * y)
\end{ReviewVerbatim}

\paragraph{Recorded source-to-library screening}
\textbf{Verdict:} \texttt{not recorded}
\reviewmatch{Matches source input}
\noindent\textbf{Reviewer annotation}\par
\reviewerbox
\clearpage
\hypertarget{library-prerequisite-aa3d377dcc497061}{}
\subsection*{\small\ttfamily\raggedright EconCSLib.\allowbreak{}Statistics.\allowbreak{}JensenConvex}
\textbf{Source connection:} no explicit byte-pinned paper source connection is registered.
\paragraph{Lean-expanded library semantic target}
\begin{ReviewVerbatim}
∀ (f : ℝ → ℝ) (x y lam : ℝ), 0 ≤ lam → lam ≤ 1 → f (lam * x + (1 - lam) * y) ≤ lam * f x + (1 - lam) * f y
\end{ReviewVerbatim}

\paragraph{Exact Lean library declaration}
\begin{ReviewVerbatim}
def JensenConvex (f : ℝ → ℝ) : Prop :=
  ∀ x y lam, 0 ≤ lam → lam ≤ 1 →
    f (lam * x + (1 - lam) * y) ≤ lam * f x + (1 - lam) * f y
\end{ReviewVerbatim}

\paragraph{Recorded source-to-library screening}
\textbf{Verdict:} \texttt{not recorded}
\reviewmatch{Matches source input}
\noindent\textbf{Reviewer annotation}\par
\reviewerbox
\clearpage
\hypertarget{library-prerequisite-e5c7eedc40c11b7a}{}
\subsection*{\small\ttfamily\raggedright EconCSLib.\allowbreak{}pmfExp}
\textbf{Source connection:} no explicit byte-pinned paper source connection is registered.
\paragraph{Lean-expanded library semantic target}
\begin{ReviewVerbatim}
{α : Type u_1} → [inst : Fintype α] → [DecidableEq α] → (μ : PMF α) → (f : α → ℝ) → ∑ a, (μ a).toReal * f a
\end{ReviewVerbatim}

\paragraph{Exact Lean library declaration}
\begin{ReviewVerbatim}
noncomputable def pmfExp {α : Type*} [Fintype α] [DecidableEq α]
    (μ : PMF α) (f : α → ℝ) : ℝ :=
  ∑ a, (μ a).toReal * f a
\end{ReviewVerbatim}

\paragraph{Recorded source-to-library screening}
\textbf{Verdict:} \texttt{not recorded}
\reviewmatch{Matches source input}
\noindent\textbf{Reviewer annotation}\par
\reviewerbox

\clearpage
\hypertarget{source-claim-7a45c0b2fffc539f}{}
\hypertarget{source-section-f10b5f68e4acf3dc}{}
\section*{Source claims}
\subsection*{Review row 1}
\textbf{Source locator:} {\footnotesize\raggedright cited publication, lines 203-\allowbreak{}276\par}
\paragraph{Verbatim source input}
\begin{ReviewVerbatim}
Beta-Bernoulli setting, this posterior distribution for product
v at timestep t is a Beta distribution with parameters
α = α̂ + R(v, t),
β = β̂ + (|S(v, t)| − R(v, t)),
where R(v, t) ∼ Binomial(t, qv ) is the number of positive
ratings received by the item so far (and so |S(v, t)| − R(v, t)
is the number of negative ratings). We can compute the posterior mean estimated quality for product v at timestep t as:

Model Primitives

Our modeling framework is as follows.
Products. There is a universe of products V . Product v ∈
V has unobserved true quality qv , where 0 ≤ qv ≤ 1.
Binary ratings that accumulate over time. Products accumulate ratings over time, as users interact with and rate
products. To capture this phenomenon, we assume the market operates over discrete time periods t, and products can
receive ratings at each time period. We assume these ratings
are binary, i.e., 0 or 1.
Crucially, the true quality of a product should impact the
rating received; we suppose that a rating for product v is
Bernoulli(qv ), independent of all other randomness in the
system. Since ratings accumulate, we use S(v, T ) to denote
the set of binary ratings that product v ∈ V has received
after T timesteps have elapsed.
We assume binary ratings as a simplification for ease of
exposition; our results qualitatively extend to other settings,
as we demonstrate in Appendix E. We note that this choice
captures how many platforms elicit ratings (e.g., in any platform that uses a thumbs-up/thumbs-down system); in other
platforms with more than two options, ratings often follow a
“J-curve” where users mostly only use the extreme options
(Hu, Zhang, and Pavlou 2009). In Appendix G, we demonstrate that our results hold even when we assume users’ distribution over ratings is centered or even right-skewed.
Prior-Weighted Rating System. In a prior-weighted rating system, the prior distribution is a design choice of the
platform. The system operates by estimating a posterior distribution of the estimated true quality for each product, given
the prior and the received ratings on that product.
We primarily consider rating systems where the prior represents a Beta distribution.1 Then, in this setting, the platform chooses the prior Beta distribution parameters (α̂, β̂).2
We will refer to α̂ + β̂ as the prior strength, as it indicates
how strongly the prior is concentrated around its mean (and
how many ratings will be needed to “overwhelm”) the prior.
After a product receives a rating, the platform updates its
estimation posterior distribution for its true quality. In our
1
This choice is natural, given that ratings are assumed to be
Bernoulli. The Beta distribution is the conjugate prior for Bernoulli
data, i.e., with this choice of prior, the posterior distribution remains a Beta distribution but with different parameters.
2
α
In a Beta(α, β) distribution, the mean is α+β
; the higher that
α + β are, the more concentrated the distribution around the mean.
A common interpretation of the parameters is that α+β is the number of trials (coin flips), and α represents the number of successes.

q̂α̂,β̂ (v, t) =

α̂ + R(v, t)
α̂ + β̂ + |S(v, t)|

.

(1)

The platform impacts the operation of a prior-weighted
rating system only through the prior choice, determined here
by α̂ and β̂. There are two aspects of this choice: the shape
of the prior, and its strength relative to observed ratings in
determining the posterior. In our analysis, we fix the prior
shape by relying on market-level data, in a manner inspired
by empirical Bayesian statistical methods (Robbins 1964).
Our analysis shows that the strength of the prior then mediates the tradeoff between efficiency and producer fairness.
Formally, to reason about a prior’s strength, we fix a prior
shape (α̃, β̃), then set (α̂, β̂) = (η α̃, η β̃) for some prior
strength η ≥ 0, providing us with a single-parameter lever
to adjust the strength of the prior. We elaborate on how we
calibrate the prior shape in Section 3.2. Note in particular
that if η = 0, then the platform simply implements sample
averaging to compute the estimated quality of an item.
\end{ReviewVerbatim}


\paragraph{Expanded PaperInterface specification}
\begin{ReviewVerbatim}
∀ (alphaHat betaHat : ℝ) (positiveRatings totalRatings : ℕ),
  MBJG25ProducerFairness.paper_posterior_rating alphaHat betaHat positiveRatings totalRatings =
    (alphaHat + ↑positiveRatings) / (alphaHat + betaHat + ↑totalRatings)
\end{ReviewVerbatim}

\paragraph{Lean proof endpoint (built separately)}\begin{ReviewVerbatim}
MBJG25ProducerFairness.paper_posterior_rating
\end{ReviewVerbatim}

\paragraph{Recorded source-to-Spec screening}
\textbf{Verdict:} \texttt{not recorded}
\\
\textbf{Reason:} not recorded
\reviewmatch{Matches source input}
\noindent\textbf{Reviewer annotation}\par
\reviewerbox
\clearpage
\hypertarget{source-claim-ec7384519273acb4}{}
\section*{Review row 2}
\textbf{Source locator:} {\footnotesize\raggedright cited publication, lines 203-\allowbreak{}276\par}
\paragraph{Verbatim source input}
\begin{ReviewVerbatim}
Beta-Bernoulli setting, this posterior distribution for product
v at timestep t is a Beta distribution with parameters
α = α̂ + R(v, t),
β = β̂ + (|S(v, t)| − R(v, t)),
where R(v, t) ∼ Binomial(t, qv ) is the number of positive
ratings received by the item so far (and so |S(v, t)| − R(v, t)
is the number of negative ratings). We can compute the posterior mean estimated quality for product v at timestep t as:

Model Primitives

Our modeling framework is as follows.
Products. There is a universe of products V . Product v ∈
V has unobserved true quality qv , where 0 ≤ qv ≤ 1.
Binary ratings that accumulate over time. Products accumulate ratings over time, as users interact with and rate
products. To capture this phenomenon, we assume the market operates over discrete time periods t, and products can
receive ratings at each time period. We assume these ratings
are binary, i.e., 0 or 1.
Crucially, the true quality of a product should impact the
rating received; we suppose that a rating for product v is
Bernoulli(qv ), independent of all other randomness in the
system. Since ratings accumulate, we use S(v, T ) to denote
the set of binary ratings that product v ∈ V has received
after T timesteps have elapsed.
We assume binary ratings as a simplification for ease of
exposition; our results qualitatively extend to other settings,
as we demonstrate in Appendix E. We note that this choice
captures how many platforms elicit ratings (e.g., in any platform that uses a thumbs-up/thumbs-down system); in other
platforms with more than two options, ratings often follow a
“J-curve” where users mostly only use the extreme options
(Hu, Zhang, and Pavlou 2009). In Appendix G, we demonstrate that our results hold even when we assume users’ distribution over ratings is centered or even right-skewed.
Prior-Weighted Rating System. In a prior-weighted rating system, the prior distribution is a design choice of the
platform. The system operates by estimating a posterior distribution of the estimated true quality for each product, given
the prior and the received ratings on that product.
We primarily consider rating systems where the prior represents a Beta distribution.1 Then, in this setting, the platform chooses the prior Beta distribution parameters (α̂, β̂).2
We will refer to α̂ + β̂ as the prior strength, as it indicates
how strongly the prior is concentrated around its mean (and
how many ratings will be needed to “overwhelm”) the prior.
After a product receives a rating, the platform updates its
estimation posterior distribution for its true quality. In our
1
This choice is natural, given that ratings are assumed to be
Bernoulli. The Beta distribution is the conjugate prior for Bernoulli
data, i.e., with this choice of prior, the posterior distribution remains a Beta distribution but with different parameters.
2
α
In a Beta(α, β) distribution, the mean is α+β
; the higher that
α + β are, the more concentrated the distribution around the mean.
A common interpretation of the parameters is that α+β is the number of trials (coin flips), and α represents the number of successes.

q̂α̂,β̂ (v, t) =

α̂ + R(v, t)
α̂ + β̂ + |S(v, t)|

.

(1)

The platform impacts the operation of a prior-weighted
rating system only through the prior choice, determined here
by α̂ and β̂. There are two aspects of this choice: the shape
of the prior, and its strength relative to observed ratings in
determining the posterior. In our analysis, we fix the prior
shape by relying on market-level data, in a manner inspired
by empirical Bayesian statistical methods (Robbins 1964).
Our analysis shows that the strength of the prior then mediates the tradeoff between efficiency and producer fairness.
Formally, to reason about a prior’s strength, we fix a prior
shape (α̃, β̃), then set (α̂, β̂) = (η α̃, η β̃) for some prior
strength η ≥ 0, providing us with a single-parameter lever
to adjust the strength of the prior. We elaborate on how we
calibrate the prior shape in Section 3.2. Note in particular
that if η = 0, then the platform simply implements sample
averaging to compute the estimated quality of an item.
\end{ReviewVerbatim}


\paragraph{Expanded PaperInterface specification}
\begin{ReviewVerbatim}
∀ (alpha beta eta t q_v : ℝ),
  MBJG25ProducerFairness.paper_posterior_mean alpha beta eta t q_v =
    (eta * alpha + t * q_v) / (eta * alpha + eta * beta + t)
\end{ReviewVerbatim}

\paragraph{Lean proof endpoint (built separately)}\begin{ReviewVerbatim}
MBJG25ProducerFairness.paper_posterior_mean
\end{ReviewVerbatim}

\paragraph{Recorded source-to-Spec screening}
\textbf{Verdict:} \texttt{not recorded}
\\
\textbf{Reason:} not recorded
\reviewmatch{Matches source input}
\noindent\textbf{Reviewer annotation}\par
\reviewerbox
\clearpage
\hypertarget{source-claim-20150c21c10be6e2}{}
\section*{Review row 3}
\textbf{Source locator:} {\footnotesize\raggedright cited publication, lines 333-\allowbreak{}341\par}
\paragraph{Verbatim source input}
\begin{ReviewVerbatim}
Why MSE? Via the bias-variance decomposition, meansquared error encodes a metric for efficiency and a metric
for producer unfairness. For product v with true quality qv :
[U+0010 control character]
[U+0011 control character]2
E[(q̂ηα̃,ηβ̃ (v, t) − qv )2 |qv ] = E[q̂ηα̃,ηβ̃ (v, t)|qv ] − qv
+ Var[q̂ηα̃,ηβ̃ (v, t)|qv ].
(3)
When conditioned on a product with given true quality
qv , the variance component of the MSE (second component
\end{ReviewVerbatim}


\paragraph{Expanded PaperInterface specification}
\begin{ReviewVerbatim}
∀ (alpha beta eta t q_v : ℝ),
  MBJG25ProducerFairness.paper_bias alpha beta eta t q_v =
    (eta * alpha + t * q_v) / (eta * alpha + eta * beta + t) - q_v
\end{ReviewVerbatim}

\paragraph{Lean proof endpoint (built separately)}\begin{ReviewVerbatim}
MBJG25ProducerFairness.paper_bias
\end{ReviewVerbatim}

\paragraph{Recorded source-to-Spec screening}
\textbf{Verdict:} \texttt{not recorded}
\\
\textbf{Reason:} not recorded
\reviewmatch{Matches source input}
\noindent\textbf{Reviewer annotation}\par
\reviewerbox
\clearpage
\hypertarget{source-claim-23bc13a84ad0b29a}{}
\section*{Review row 4}
\textbf{Source locator:} {\footnotesize\raggedright cited publication, lines 1533-\allowbreak{}1552\par}
\paragraph{Verbatim source input}
\begin{ReviewVerbatim}
Variance can also be broken down accordingly:
"
Var[q̂ηα̃,ηβ̃ (v, t)|qv ] = E

[U+0010 control character]

[U+0011 control character]2
q̂ηα̃,ηβ̃ (v, t) − E[q̂ηα̃,ηβ̃ (v, t)] qv
(8)
2

E[(R(v, t) − E[R(v, t)]) |qv ]
(η α̃ + η β̃ + t)2
Var[R(v, t)|qv ]
=
(η α̃ + η β̃ + t)2
tqv (1 − qv )
=
.
(η α̃ + η β̃ + t)2
\end{ReviewVerbatim}


\paragraph{Expanded PaperInterface specification}
\begin{ReviewVerbatim}
∀ (alpha beta eta t q_v : ℝ),
  MBJG25ProducerFairness.paper_variance alpha beta eta t q_v = t * q_v * (1 - q_v) / (eta * alpha + eta * beta + t) ^ 2
\end{ReviewVerbatim}

\paragraph{Lean proof endpoint (built separately)}\begin{ReviewVerbatim}
MBJG25ProducerFairness.paper_variance
\end{ReviewVerbatim}

\paragraph{Recorded source-to-Spec screening}
\textbf{Verdict:} \texttt{not recorded}
\\
\textbf{Reason:} not recorded
\reviewmatch{Matches source input}
\noindent\textbf{Reviewer annotation}\par
\reviewerbox
\clearpage
\hypertarget{source-claim-6e52f3be6106ad49}{}
\section*{Review row 5}
\textbf{Source locator:} {\footnotesize\raggedright cited publication, lines 333-\allowbreak{}341\par}
\paragraph{Verbatim source input}
\begin{ReviewVerbatim}
Why MSE? Via the bias-variance decomposition, meansquared error encodes a metric for efficiency and a metric
for producer unfairness. For product v with true quality qv :
[U+0010 control character]
[U+0011 control character]2
E[(q̂ηα̃,ηβ̃ (v, t) − qv )2 |qv ] = E[q̂ηα̃,ηβ̃ (v, t)|qv ] − qv
+ Var[q̂ηα̃,ηβ̃ (v, t)|qv ].
(3)
When conditioned on a product with given true quality
qv , the variance component of the MSE (second component
\end{ReviewVerbatim}


\paragraph{Expanded PaperInterface specification}
\begin{ReviewVerbatim}
∀ (alpha beta eta t q_v : ℝ),
  MBJG25ProducerFairness.paper_squared_bias alpha beta eta t q_v =
    MBJG25ProducerFairness.paper_bias alpha beta eta t q_v ^ 2
\end{ReviewVerbatim}

\paragraph{Lean proof endpoint (built separately)}\begin{ReviewVerbatim}
MBJG25ProducerFairness.paper_squared_bias
\end{ReviewVerbatim}

\paragraph{Recorded source-to-Spec screening}
\textbf{Verdict:} \texttt{not recorded}
\\
\textbf{Reason:} not recorded
\reviewmatch{Matches source input}
\noindent\textbf{Reviewer annotation}\par
\reviewerbox
\clearpage
\hypertarget{source-claim-9fd08233a5aeb857}{}
\section*{Review row 6}
\textbf{Source locator:} {\footnotesize\raggedright cited publication, lines 315-\allowbreak{}333\par}
\paragraph{Verbatim source input}
\begin{ReviewVerbatim}
Informally, efficiency is measured by whether or not the rating system is successfully identifying high-quality products.
Individual producer fairness is measured by whether products of similar true quality are treated similarly.
Metrics in fixed setting. In the fixed model, we focus
primarily on an accuracy measure: the mean-squared error
(MSE) of the estimated quality in predicting true quality,
defined at timestep t as
P
M SE(q̂ηα̃,ηβ̃ ) =

2
v∈V (q̂η α̃,η β̃ (v, t) − qv )

|V |

.

(2)

Why MSE? Via the bias-variance decomposition, meansquared error encodes a metric for efficiency and a metric
\end{ReviewVerbatim}


\paragraph{Expanded PaperInterface specification}
\begin{ReviewVerbatim}
∀ {V : Type u_1} [inst : Fintype V] [inst_1 : Nonempty V] (estimatedQuality trueQuality : V → ℝ),
  MBJG25ProducerFairness.paper_fixed_market_mse estimatedQuality trueQuality =
    (∑ v, (estimatedQuality v - trueQuality v) ^ 2) / ↑(Fintype.card V)
\end{ReviewVerbatim}

\paragraph{Lean proof endpoint (built separately)}\begin{ReviewVerbatim}
MBJG25ProducerFairness.paper_fixed_market_mse
\end{ReviewVerbatim}

\paragraph{Recorded source-to-Spec screening}
\textbf{Verdict:} \texttt{not recorded}
\\
\textbf{Reason:} not recorded
\reviewmatch{Matches source input}
\noindent\textbf{Reviewer annotation}\par
\reviewerbox
\clearpage
\hypertarget{source-claim-da6d1b14c0df1e3e}{}
\section*{Review row 7}
\textbf{Source locator:} {\footnotesize\raggedright cited publication, lines 333-\allowbreak{}345\par}
\paragraph{Verbatim source input}
\begin{ReviewVerbatim}
Why MSE? Via the bias-variance decomposition, meansquared error encodes a metric for efficiency and a metric
for producer unfairness. For product v with true quality qv :
[U+0010 control character]
[U+0011 control character]2
E[(q̂ηα̃,ηβ̃ (v, t) − qv )2 |qv ] = E[q̂ηα̃,ηβ̃ (v, t)|qv ] − qv
+ Var[q̂ηα̃,ηβ̃ (v, t)|qv ].
(3)
When conditioned on a product with given true quality
qv , the variance component of the MSE (second component
on the right-hand side) represents an objective for producer
fairness, as the estimator should ideally estimate products
of similar true quality similarly, even if that estimated quality is biased. The bias component of MSE is correlated with
consumer efficiency; in particular, a rating system that consistently underrates high-quality products or overrates lowquality ones will lead to poor consumer experiences.
\end{ReviewVerbatim}


\paragraph{Expanded PaperInterface specification}
\begin{ReviewVerbatim}
∀ {Ω : Type u_1} [inst : Fintype Ω] [inst_1 : DecidableEq Ω] (rating_dist : PMF Ω) (posterior_rating : Ω → ℝ) (q_v : ℝ),
  (EconCSLib.pmfExp rating_dist fun ω => (posterior_rating ω - q_v) ^ 2) =
    (EconCSLib.pmfExp rating_dist posterior_rating - q_v) ^ 2 +
      EconCSLib.pmfExp rating_dist fun ω => (posterior_rating ω - EconCSLib.pmfExp rating_dist posterior_rating) ^ 2
\end{ReviewVerbatim}

\paragraph{Lean proof endpoint (built separately)}\begin{ReviewVerbatim}
MBJG25ProducerFairness.paper_facing_fixed_mse_decomposition
\end{ReviewVerbatim}

\paragraph{Recorded source-to-Spec screening}
\textbf{Verdict:} \texttt{not recorded}
\\
\textbf{Reason:} not recorded
\reviewmatch{Matches source input}
\noindent\textbf{Reviewer annotation}\par
\reviewerbox
\clearpage
\hypertarget{source-claim-000171f0bf87b98c}{}
\section*{Review row 8}
\textbf{Source locator:} {\footnotesize\raggedright cited publication, lines 443-\allowbreak{}450\par}
\paragraph{Verbatim source input}
\begin{ReviewVerbatim}
Theorem 3.1. Suppose we have a prior-weighted rating system in the fixed setting with prior parameters (η α̃, η β̃). Fix
product v with true quality qv and consider quality estimation after t timesteps.
Then, as η
≥ 0 increases, (1) squared bias
[U+0010 control character]
[U+0011 control character]2
E[q̂ηα̃,ηβ̃ (v, t)|qv ] − qv
is nondecreasing; (2) variance Var[q̂ηα̃,ηβ̃ (v, t)|qv ] is strictly decreasing in η.
\end{ReviewVerbatim}


\paragraph{Expanded PaperInterface specification}
\begin{ReviewVerbatim}
∀ {alpha beta t q etaLow etaHigh : ℝ},
  MBJG25ProducerFairness.assumption_positive_prior_shape alpha beta →
    MBJG25ProducerFairness.assumption_positive_time t →
      MBJG25ProducerFairness.assumption_quality_nonnegative q →
        MBJG25ProducerFairness.assumption_quality_at_most_one q →
          MBJG25ProducerFairness.assumption_prior_strength_nonnegative etaLow →
            MBJG25ProducerFairness.assumption_prior_strength_weak_order etaLow etaHigh →
              MBJG25ProducerFairness.paper_variance alpha beta etaHigh t q ≤
                MBJG25ProducerFairness.paper_variance alpha beta etaLow t q
\end{ReviewVerbatim}

\paragraph{Lean proof endpoint (built separately)}\begin{ReviewVerbatim}
MBJG25ProducerFairness.paper_facing_theorem3_1_variance_weak_decrease
\end{ReviewVerbatim}

\paragraph{Recorded source-to-Spec screening}
\textbf{Verdict:} \texttt{not recorded}
\\
\textbf{Reason:} not recorded
\reviewmatch{Matches source input}
\noindent\textbf{Reviewer annotation}\par
\reviewerbox
\clearpage
\hypertarget{source-claim-976fc5a9912d46b4}{}
\section*{Review row 9}
\textbf{Source locator:} {\footnotesize\raggedright cited publication, lines 443-\allowbreak{}450\par}
\paragraph{Verbatim source input}
\begin{ReviewVerbatim}
Theorem 3.1. Suppose we have a prior-weighted rating system in the fixed setting with prior parameters (η α̃, η β̃). Fix
product v with true quality qv and consider quality estimation after t timesteps.
Then, as η
≥ 0 increases, (1) squared bias
[U+0010 control character]
[U+0011 control character]2
E[q̂ηα̃,ηβ̃ (v, t)|qv ] − qv
is nondecreasing; (2) variance Var[q̂ηα̃,ηβ̃ (v, t)|qv ] is strictly decreasing in η.
\end{ReviewVerbatim}


\paragraph{Expanded PaperInterface specification}
\begin{ReviewVerbatim}
∀ {alpha beta t q etaLow etaHigh : ℝ},
  MBJG25ProducerFairness.assumption_positive_prior_shape alpha beta →
    MBJG25ProducerFairness.assumption_positive_time t →
      MBJG25ProducerFairness.assumption_quality_positive q →
        MBJG25ProducerFairness.assumption_quality_lt_one q →
          MBJG25ProducerFairness.assumption_prior_strength_nonnegative etaLow →
            MBJG25ProducerFairness.assumption_prior_strength_strict_order etaLow etaHigh →
              MBJG25ProducerFairness.paper_variance alpha beta etaHigh t q <
                MBJG25ProducerFairness.paper_variance alpha beta etaLow t q
\end{ReviewVerbatim}

\paragraph{Lean proof endpoint (built separately)}\begin{ReviewVerbatim}
MBJG25ProducerFairness.paper_facing_theorem3_1_variance_strict_decrease_interior
\end{ReviewVerbatim}

\paragraph{Recorded source-to-Spec screening}
\textbf{Verdict:} \texttt{not recorded}
\\
\textbf{Reason:} not recorded
\reviewmatch{Matches source input}
\noindent\textbf{Reviewer annotation}\par
\reviewerbox
\clearpage
\hypertarget{source-claim-5564d9d725c13f20}{}
\section*{Review row 10}
\textbf{Source locator:} {\footnotesize\raggedright cited publication, lines 443-\allowbreak{}450\par}
\paragraph{Verbatim source input}
\begin{ReviewVerbatim}
Theorem 3.1. Suppose we have a prior-weighted rating system in the fixed setting with prior parameters (η α̃, η β̃). Fix
product v with true quality qv and consider quality estimation after t timesteps.
Then, as η
≥ 0 increases, (1) squared bias
[U+0010 control character]
[U+0011 control character]2
E[q̂ηα̃,ηβ̃ (v, t)|qv ] − qv
is nondecreasing; (2) variance Var[q̂ηα̃,ηβ̃ (v, t)|qv ] is strictly decreasing in η.
\end{ReviewVerbatim}


\paragraph{Expanded PaperInterface specification}
\begin{ReviewVerbatim}
∀ {alpha beta t q etaLow etaHigh : ℝ},
  MBJG25ProducerFairness.assumption_positive_prior_shape alpha beta →
    MBJG25ProducerFairness.assumption_positive_time t →
      MBJG25ProducerFairness.assumption_prior_strength_nonnegative etaLow →
        MBJG25ProducerFairness.assumption_prior_strength_weak_order etaLow etaHigh →
          MBJG25ProducerFairness.paper_squared_bias alpha beta etaLow t q ≤
            MBJG25ProducerFairness.paper_squared_bias alpha beta etaHigh t q
\end{ReviewVerbatim}

\paragraph{Lean proof endpoint (built separately)}\begin{ReviewVerbatim}
MBJG25ProducerFairness.paper_facing_theorem3_1_squared_bias_nondecreasing
\end{ReviewVerbatim}

\paragraph{Recorded source-to-Spec screening}
\textbf{Verdict:} \texttt{not recorded}
\\
\textbf{Reason:} not recorded
\reviewmatch{Matches source input}
\noindent\textbf{Reviewer annotation}\par
\reviewerbox
\clearpage
\hypertarget{source-claim-2dcf6dbcf162b5d0}{}
\section*{Review row 11}
\textbf{Source locator:} {\footnotesize\raggedright cited publication, lines 460-\allowbreak{}469\par}
\paragraph{Verbatim source input}
\begin{ReviewVerbatim}
Theorem 3.2. Suppose we have a prior-weighted rating system in the fixed setting with prior parameters (η α̃, η β̃). Consider quality estimation after t timesteps.
Then, as a function of true product quality, squared bias
[U+0010 control character]
[U+0011 control character]2
E[q̂ηα̃,ηβ̃ (v, t)|qv ] − qv
is convex with a global miniα̃
mum at α̃+
, while variance Var[q̂ηα̃,ηβ̃ (v, t)|qv ] is concave
β̃
with a global maximum at 1/2.
\end{ReviewVerbatim}


\paragraph{Expanded PaperInterface specification}
\begin{ReviewVerbatim}
∀ {alpha beta eta t : ℝ},
  MBJG25ProducerFairness.assumption_positive_prior_shape alpha beta →
    MBJG25ProducerFairness.assumption_prior_strength_nonnegative eta →
      MBJG25ProducerFairness.assumption_positive_time t →
        EconCSLib.Statistics.JensenConvex fun q => MBJG25ProducerFairness.paper_squared_bias alpha beta eta t q
\end{ReviewVerbatim}

\paragraph{Lean proof endpoint (built separately)}\begin{ReviewVerbatim}
MBJG25ProducerFairness.paper_facing_theorem3_2_squared_bias_convex_in_quality
\end{ReviewVerbatim}

\paragraph{Recorded source-to-Spec screening}
\textbf{Verdict:} \texttt{not recorded}
\\
\textbf{Reason:} not recorded
\reviewmatch{Matches source input}
\noindent\textbf{Reviewer annotation}\par
\reviewerbox
\clearpage
\hypertarget{source-claim-444ef89e21ba4476}{}
\section*{Review row 12}
\textbf{Source locator:} {\footnotesize\raggedright cited publication, lines 460-\allowbreak{}469\par}
\paragraph{Verbatim source input}
\begin{ReviewVerbatim}
Theorem 3.2. Suppose we have a prior-weighted rating system in the fixed setting with prior parameters (η α̃, η β̃). Consider quality estimation after t timesteps.
Then, as a function of true product quality, squared bias
[U+0010 control character]
[U+0011 control character]2
E[q̂ηα̃,ηβ̃ (v, t)|qv ] − qv
is convex with a global miniα̃
mum at α̃+
, while variance Var[q̂ηα̃,ηβ̃ (v, t)|qv ] is concave
β̃
with a global maximum at 1/2.
\end{ReviewVerbatim}


\paragraph{Expanded PaperInterface specification}
\begin{ReviewVerbatim}
∀ {alpha beta eta t : ℝ},
  MBJG25ProducerFairness.assumption_positive_prior_shape alpha beta →
    MBJG25ProducerFairness.assumption_prior_strength_nonnegative eta →
      MBJG25ProducerFairness.assumption_positive_time t →
        EconCSLib.Statistics.GlobalMinAt (fun q => MBJG25ProducerFairness.paper_squared_bias alpha beta eta t q)
          (alpha / (alpha + beta))
\end{ReviewVerbatim}

\paragraph{Lean proof endpoint (built separately)}\begin{ReviewVerbatim}
MBJG25ProducerFairness.paper_facing_theorem3_2_squared_bias_global_min_at_prior_mean
\end{ReviewVerbatim}

\paragraph{Recorded source-to-Spec screening}
\textbf{Verdict:} \texttt{not recorded}
\\
\textbf{Reason:} not recorded
\reviewmatch{Matches source input}
\noindent\textbf{Reviewer annotation}\par
\reviewerbox
\clearpage
\hypertarget{source-claim-3afe5de598f0ed9d}{}
\section*{Review row 13}
\textbf{Source locator:} {\footnotesize\raggedright cited publication, lines 460-\allowbreak{}469\par}
\paragraph{Verbatim source input}
\begin{ReviewVerbatim}
Theorem 3.2. Suppose we have a prior-weighted rating system in the fixed setting with prior parameters (η α̃, η β̃). Consider quality estimation after t timesteps.
Then, as a function of true product quality, squared bias
[U+0010 control character]
[U+0011 control character]2
E[q̂ηα̃,ηβ̃ (v, t)|qv ] − qv
is convex with a global miniα̃
mum at α̃+
, while variance Var[q̂ηα̃,ηβ̃ (v, t)|qv ] is concave
β̃
with a global maximum at 1/2.
\end{ReviewVerbatim}


\paragraph{Expanded PaperInterface specification}
\begin{ReviewVerbatim}
∀ {alpha beta eta t : ℝ},
  MBJG25ProducerFairness.assumption_nonnegative_time t →
    EconCSLib.Statistics.JensenConcave fun q => MBJG25ProducerFairness.paper_variance alpha beta eta t q
\end{ReviewVerbatim}

\paragraph{Lean proof endpoint (built separately)}\begin{ReviewVerbatim}
MBJG25ProducerFairness.paper_facing_theorem3_2_variance_concave_in_quality
\end{ReviewVerbatim}

\paragraph{Recorded source-to-Spec screening}
\textbf{Verdict:} \texttt{not recorded}
\\
\textbf{Reason:} not recorded
\reviewmatch{Matches source input}
\noindent\textbf{Reviewer annotation}\par
\reviewerbox
\clearpage
\hypertarget{source-claim-a4656e58047f308c}{}
\section*{Review row 14}
\textbf{Source locator:} {\footnotesize\raggedright cited publication, lines 460-\allowbreak{}469\par}
\paragraph{Verbatim source input}
\begin{ReviewVerbatim}
Theorem 3.2. Suppose we have a prior-weighted rating system in the fixed setting with prior parameters (η α̃, η β̃). Consider quality estimation after t timesteps.
Then, as a function of true product quality, squared bias
[U+0010 control character]
[U+0011 control character]2
E[q̂ηα̃,ηβ̃ (v, t)|qv ] − qv
is convex with a global miniα̃
mum at α̃+
, while variance Var[q̂ηα̃,ηβ̃ (v, t)|qv ] is concave
β̃
with a global maximum at 1/2.
\end{ReviewVerbatim}


\paragraph{Expanded PaperInterface specification}
\begin{ReviewVerbatim}
∀ {alpha beta eta t : ℝ},
  MBJG25ProducerFairness.assumption_nonnegative_time t →
    EconCSLib.Statistics.GlobalMaxAt (fun q => MBJG25ProducerFairness.paper_variance alpha beta eta t q) (1 / 2)
\end{ReviewVerbatim}

\paragraph{Lean proof endpoint (built separately)}\begin{ReviewVerbatim}
MBJG25ProducerFairness.paper_facing_theorem3_2_variance_global_max_at_half
\end{ReviewVerbatim}

\paragraph{Recorded source-to-Spec screening}
\textbf{Verdict:} \texttt{not recorded}
\\
\textbf{Reason:} not recorded
\reviewmatch{Matches source input}
\noindent\textbf{Reviewer annotation}\par
\reviewerbox
\clearpage
\hypertarget{source-claim-6ce0e529b89e57b7}{}
\section*{Review row 15}
\textbf{Source locator:} {\footnotesize\raggedright cited publication, lines 443-\allowbreak{}450\par}
\paragraph{Verbatim source input}
\begin{ReviewVerbatim}
Theorem 3.1. Suppose we have a prior-weighted rating system in the fixed setting with prior parameters (η α̃, η β̃). Fix
product v with true quality qv and consider quality estimation after t timesteps.
Then, as η
≥ 0 increases, (1) squared bias
[U+0010 control character]
[U+0011 control character]2
E[q̂ηα̃,ηβ̃ (v, t)|qv ] − qv
is nondecreasing; (2) variance Var[q̂ηα̃,ηβ̃ (v, t)|qv ] is strictly decreasing in η.
\end{ReviewVerbatim}


\paragraph{Expanded PaperInterface specification}
\begin{ReviewVerbatim}
∀ (alpha beta t etaLow etaHigh : ℝ),
  ¬MBJG25ProducerFairness.paper_variance alpha beta etaHigh t 0 <
      MBJG25ProducerFairness.paper_variance alpha beta etaLow t 0
\end{ReviewVerbatim}

\paragraph{Lean proof endpoint (built separately)}\begin{ReviewVerbatim}
MBJG25ProducerFairness.paper_facing_theorem3_1_variance_strict_decrease_counterexample_quality_zero
\end{ReviewVerbatim}

\paragraph{Recorded source-to-Spec screening}
\textbf{Verdict:} \texttt{not recorded}
\\
\textbf{Reason:} not recorded
\reviewmatch{Matches source input}
\noindent\textbf{Reviewer annotation}\par
\reviewerbox
\clearpage
\hypertarget{source-claim-464dc3b105b500a1}{}
\section*{Review row 16}
\textbf{Source locator:} {\footnotesize\raggedright cited publication, lines 443-\allowbreak{}450\par}
\paragraph{Verbatim source input}
\begin{ReviewVerbatim}
Theorem 3.1. Suppose we have a prior-weighted rating system in the fixed setting with prior parameters (η α̃, η β̃). Fix
product v with true quality qv and consider quality estimation after t timesteps.
Then, as η
≥ 0 increases, (1) squared bias
[U+0010 control character]
[U+0011 control character]2
E[q̂ηα̃,ηβ̃ (v, t)|qv ] − qv
is nondecreasing; (2) variance Var[q̂ηα̃,ηβ̃ (v, t)|qv ] is strictly decreasing in η.
\end{ReviewVerbatim}


\paragraph{Expanded PaperInterface specification}
\begin{ReviewVerbatim}
∀ (alpha beta t etaLow etaHigh : ℝ),
  ¬MBJG25ProducerFairness.paper_variance alpha beta etaHigh t 1 <
      MBJG25ProducerFairness.paper_variance alpha beta etaLow t 1
\end{ReviewVerbatim}

\paragraph{Lean proof endpoint (built separately)}\begin{ReviewVerbatim}
MBJG25ProducerFairness.paper_facing_theorem3_1_variance_strict_decrease_counterexample_quality_one
\end{ReviewVerbatim}

\paragraph{Recorded source-to-Spec screening}
\textbf{Verdict:} \texttt{not recorded}
\\
\textbf{Reason:} not recorded
\reviewmatch{Matches source input}
\noindent\textbf{Reviewer annotation}\par
\reviewerbox
\clearpage
\hypertarget{source-claim-5e3bedb3cf1d62b0}{}
\section*{Review row 17}
\textbf{Source locator:} {\footnotesize\raggedright cited publication, lines 346-\allowbreak{}353\par}
\paragraph{Verbatim source input}
\begin{ReviewVerbatim}
Metrics in responsive setting. In the responsive setting,
because consumers choose (and rate) higher rated items at
higher probabilities, we base our metrics directly off of how
often products are selected. We define lv to be the lifespan
of a product v in the responsive model. The selection rate
of a product v as the fraction of times that it was rated (and
thus the fraction of times the product was purchased) during
its lifespan, SR(v) := |S(v, t)|/lv .
\end{ReviewVerbatim}


\paragraph{Expanded PaperInterface specification}
\begin{ReviewVerbatim}
∀ (selections lifespan : ℝ),
  MBJG25ProducerFairness.paper_facing_selection_rate selections lifespan = selections / lifespan
\end{ReviewVerbatim}

\paragraph{Lean proof endpoint (built separately)}\begin{ReviewVerbatim}
MBJG25ProducerFairness.paper_facing_selection_rate
\end{ReviewVerbatim}

\paragraph{Recorded source-to-Spec screening}
\textbf{Verdict:} \texttt{not recorded}
\\
\textbf{Reason:} not recorded
\reviewmatch{Matches source input}
\noindent\textbf{Reviewer annotation}\par
\reviewerbox
\clearpage
\hypertarget{source-claim-430cb3de083df7cf}{}
\section*{Review row 18}
\textbf{Source locator:} {\footnotesize\raggedright cited publication, lines 376-\allowbreak{}386\par}
\paragraph{Verbatim source input}
\begin{ReviewVerbatim}
Fairness. To capture variance conditional on quality, our
chosen fairness metric is the standard deviation in selection rate given true quality. This individual fairness notion
is conditional on quality: items of similar quality ought to
be picked at similar rates. Let σ̂(U ) be the sample standard
deviation computed on a set of real-valued numbers U . For a
fixed true quality q, the individual (un)fairness metric is the
standard deviation in selection rate conditioned on an item
being of that true quality:
σ̂ ({SR(v)|qv = q}) .

(5)
\end{ReviewVerbatim}


\paragraph{Expanded PaperInterface specification}
\begin{ReviewVerbatim}
∀ {V : Type u_1} [inst : Fintype V] [inst_1 : DecidableEq V] (selections lifespan q_v : V → ℝ) (q : ℝ),
  MBJG25ProducerFairness.paper_facing_individual_producer_unfairness selections lifespan q_v q =
    have S := {v | q_v v = q};
    have selectionRate := fun v => MBJG25ProducerFairness.paper_facing_selection_rate (selections v) (lifespan v);
    have meanSelectionRate := if S.card = 0 then 0 else (∑ v ∈ S, selectionRate v) / ↑S.card;
    if S.card = 0 then 0 else √((∑ v ∈ S, (selectionRate v - meanSelectionRate) ^ 2) / ↑S.card)
\end{ReviewVerbatim}

\paragraph{Lean proof endpoint (built separately)}\begin{ReviewVerbatim}
MBJG25ProducerFairness.paper_facing_individual_producer_unfairness
\end{ReviewVerbatim}

\paragraph{Recorded source-to-Spec screening}
\textbf{Verdict:} \texttt{not recorded}
\\
\textbf{Reason:} not recorded
\reviewmatch{Matches source input}
\noindent\textbf{Reviewer annotation}\par
\reviewerbox
\clearpage
\hypertarget{source-claim-aace03b9ed863bb2}{}
\section*{Review row 19}
\textbf{Source locator:} {\footnotesize\raggedright cited publication, lines 386-\allowbreak{}390\par}
\paragraph{Verbatim source input}
\begin{ReviewVerbatim}
(5)

To get a marketplace-level metric of producer unfairness,
we take the expectation of eq. (5) over true qualities q.
We do not use a statistical measure as in the fixed setting for two reasons: (1) The responsive model allows for
\end{ReviewVerbatim}


\paragraph{Expanded PaperInterface specification}
\begin{ReviewVerbatim}
∀ {Q : Type u_1} [inst : Fintype Q] [inst_1 : DecidableEq Q] (quality_dist : PMF Q) (quality_unfairness : Q → ℝ),
  MBJG25ProducerFairness.paper_facing_marketplace_producer_unfairness quality_dist quality_unfairness =
    ∑ q, (quality_dist q).toReal * quality_unfairness q
\end{ReviewVerbatim}

\paragraph{Lean proof endpoint (built separately)}\begin{ReviewVerbatim}
MBJG25ProducerFairness.paper_facing_marketplace_producer_unfairness
\end{ReviewVerbatim}

\paragraph{Recorded source-to-Spec screening}
\textbf{Verdict:} \texttt{not recorded}
\\
\textbf{Reason:} not recorded
\reviewmatch{Matches source input}
\noindent\textbf{Reviewer annotation}\par
\reviewerbox
\clearpage
\hypertarget{source-claim-f5087ebba0ec2e7e}{}
\section*{Review row 20}
\textbf{Source locator:} {\footnotesize\raggedright cited publication, lines 680-\allowbreak{}708\par}
\paragraph{Verbatim source input}
\begin{ReviewVerbatim}
sampling algorithm. This section is organized as follows. We

first outline how we implement simulations for the responsive model, and then investigate how tuning prior strength
affects the producer fairness and match quality metrics. We
conclude with a discussion of results, and compare our simulation outcomes with those from the fixed model.

4.1

Responsive Model Implementation

Here, we describe our simulation implementation of the responsive model. (Note that, as discussed in Section 2.3, the
changes in the responsive setting make quality estimation
inherently biased, complicating the theory developed in the
fixed setting; Appendix C demonstrates the dependency of
bias and variance on the exact dynamics of the sampling
algorithm used for consumer choice in this setting.) An
overview of the responsive model’s design and metrics for
efficiency and fairness are provided in Section 2.
Users adaptively choose products on the market. At
each timestep t, a subset of the products M (t) ⊆ V is available for purchase. A single user interacts with a single item
at each timestep, making their choice as a function of past
ratings. In particular, in the main text, we model this user as
a Thompson sampler that utilizes the posterior quality distribution for each item. That is, for each v ∈ M (t), the consumer draws a sample xv ∼ Beta(α̂+R(v, t), β̂+|S(v, t)|−
R(v, t)), and then chooses and rates the product v(t) that
maximizes xv .4 Thompson sampling is a well-studied algorithm for trading off between exploitation and exploration
in multi-armed bandit problems (Russo et al. 2018). Prior
work has also argued that it serves as a reasonable approximation of real-world consumer choice; for example, Krafft
et al. (2021) demonstrate that when individual human decisions are aggregated together in consumer financial markets, the behavior is remarkably close to that of a distributed
Thompson sampling setting.
\end{ReviewVerbatim}


\paragraph{Expanded PaperInterface specification}
\begin{ReviewVerbatim}
∀ {V : Type u_1} [Fintype V] [DecidableEq V] [Nonempty V] (belief : PMF (V → ℝ)),
  ∃ tie_breaker,
    (∀ (profile : V → ℝ) (v : V), profile v ≤ profile (tie_breaker profile)) ∧
      ∃ policy, policy = belief.bind fun profile => PMF.pure (tie_breaker profile)
\end{ReviewVerbatim}

\paragraph{Lean proof endpoint (built separately)}\begin{ReviewVerbatim}
MBJG25ProducerFairness.paper_facing_thompson_sampling_mechanism
\end{ReviewVerbatim}

\paragraph{Recorded source-to-Spec screening}
\textbf{Verdict:} \texttt{not recorded}
\\
\textbf{Reason:} not recorded
\reviewmatch{Matches source input}
\noindent\textbf{Reviewer annotation}\par
\reviewerbox
\clearpage
\hypertarget{source-claim-9acd03585d100ede}{}
\section*{Review row 21}
\textbf{Source locator:} {\footnotesize\raggedright cited publication, lines 354-\allowbreak{}369\par}
\paragraph{Verbatim source input}
\begin{ReviewVerbatim}
Efficiency. Formally, to measure efficiency, we use total
expected regret compared to the best item in the marketplace
at each time step:

E

" T
X
t=1

#
max {qv } − qv(t) .

v∈M (t)

(4)
\end{ReviewVerbatim}


\paragraph{Expanded PaperInterface specification}
\begin{ReviewVerbatim}
∀ {V : Type u_1} [inst : Fintype V] [inst_1 : DecidableEq V] [inst_2 : Nonempty V] (T : ℕ) (M : Fin T → Finset V)
  (h_nonempty : ∀ (t : Fin T), (M t).Nonempty) (q : V → ℝ) (pi : (t : Fin T) → PMF ↥(M t)),
  MBJG25ProducerFairness.paper_facing_expected_regret T M h_nonempty q pi =
    ∑ t, ((M t).sup' ⋯ q - EconCSLib.pmfExp (pi t) fun v => q ↑v)
\end{ReviewVerbatim}

\paragraph{Lean proof endpoint (built separately)}\begin{ReviewVerbatim}
MBJG25ProducerFairness.paper_facing_expected_regret
\end{ReviewVerbatim}

\paragraph{Recorded source-to-Spec screening}
\textbf{Verdict:} \texttt{not recorded}
\\
\textbf{Reason:} not recorded
\reviewmatch{Matches source input}
\noindent\textbf{Reviewer annotation}\par
\reviewerbox
\clearpage
\hypertarget{source-claim-7a7d1cda2bb14e73}{}
\section*{Review row 22}
\textbf{Source locator:} {\footnotesize\raggedright cited publication, lines 1750-\allowbreak{}1815\par}
\paragraph{Verbatim source input}
\begin{ReviewVerbatim}
50

1
Figure 8: Lineplot of mean-squared
error different quality
estimators in fixed setting over time. η = 0 represents the
sample mean estimator, η = 1 represents the empirical
Bayes (EB) estimator. The sample mean estimator (η = 0)
has high MSE for low time-steps, due to variance in ratings. On the other extreme, for extremely high prior strength
η = 1000, the platform never learns from ratings data. This
presents evidence that proper calibration of prior strength
can help solve the cold start problem in the context of accurate quality estimates for products.

Discussion of bias-variance decomposition
in responsive model

This appendix discusses a decomposition of the biasvariance decomposition in Equation 7 as applied to the responsive setting. While mean-squared prediction error may
no longer serve as an ideal metric for consumer and producer welfare in this setting, we demonstrate that markets
that care about MSE for its own sake can still reason about
prediction error for a prior-weighted rating system so long as
they can make statements about the changes in how products
are sampled based on changes in qv and η.
Let the number of reviews for product v with fixed true
quality qv , prior parameters α̃, β̃ and system-level prior
strength η, calculated at timestep t, be Nα̃,β̃ (qv , η, t). We
focus on the number of reviews for a product at a fixed
time t with fixed prior parameters α̃, β̃, and so we simplify by omitting these variables in this notation, letting
N (qv , η) denote the number of reviews. Now, at this time
t, a prior-weighted rating system’s estimated quality for v
is given by q̂ηα̃,ηβ̃ (v, N (qv , η)), and the MSE is given by
E[(q̂ηα̃,ηβ̃ (v, N (qv , η)) − qv )2 |qv ]. We can break down bias
and variance of mean-squared prediction error through the
usage of the tower law of probability:
E[(q̂ηα̃,ηβ̃ (v, N (qv , η)) − qv )2 |qv ]
h
i
= E E[(q̂ηα̃,ηβ̃ (v, N (qv , η)) − qv )2 ] N (qv , η), qv

(19)


[form-feed in source extraction]
When the number of reviews at a given point in time is
fixed, we can treat the estimated quality identically to how it
is calculated in the fixed setting. This gives us the following
bias-variance breakdown:
E[(q̂ηα̃,ηβ̃ (v, N (qv , η)) − qv )2 |qv ]
[U+F8EE private-use glyph]
!2 [U+F8F9 private-use glyph]
η α̃ − η α̃qv − η β̃qv
qv [U+F8FB private-use glyph]
= E[U+F8F0 private-use glyph]
η α̃ + η β̃ + N (qv , η)
"
#
N (qv , η)qv (1 − qv )
+E
qv
(η α̃ + η β̃ + N (qv , η))2

can be achieved by setting the prior strength η in a priorweighted rating system to zero. For a product with n ratings, k of which are positive, a population average rating of
C, and an integer-valued threshold m, the Bayesian rating q
of a product is calculated to be
n
k
m
· +
·C
\end{ReviewVerbatim}


\paragraph{Expanded PaperInterface specification}
\begin{ReviewVerbatim}
∀ {α : Type u_1} {Ω : Type u_2} [inst : Fintype α] [inst_1 : DecidableEq α] [inst_2 : Fintype Ω]
  [inst_3 : DecidableEq Ω] {alpha beta eta q_v : ℝ} (state_dist : PMF α) (rating_dist : α → PMF Ω) (N : α → ℝ)
  (posterior_rating : α → Ω → ℝ),
  (∀ (s : α),
      EconCSLib.pmfExp (rating_dist s) (posterior_rating s) =
        MBJG25ProducerFairness.paper_posterior_mean alpha beta eta (N s) q_v) →
    (∀ (s : α),
        (EconCSLib.pmfExp (rating_dist s) fun ω =>
            (posterior_rating s ω - EconCSLib.pmfExp (rating_dist s) (posterior_rating s)) ^ 2) =
          MBJG25ProducerFairness.paper_variance alpha beta eta (N s) q_v) →
      (EconCSLib.pmfExp state_dist fun s =>
          EconCSLib.pmfExp (rating_dist s) fun ω => (posterior_rating s ω - q_v) ^ 2) =
        (EconCSLib.pmfExp state_dist fun s => MBJG25ProducerFairness.paper_squared_bias alpha beta eta (N s) q_v) +
          EconCSLib.pmfExp state_dist fun s => MBJG25ProducerFairness.paper_variance alpha beta eta (N s) q_v
\end{ReviewVerbatim}

\paragraph{Lean proof endpoint (built separately)}\begin{ReviewVerbatim}
MBJG25ProducerFairness.paper_facing_responsive_mse_decomposition
\end{ReviewVerbatim}

\paragraph{Recorded source-to-Spec screening}
\textbf{Verdict:} \texttt{not recorded}
\\
\textbf{Reason:} not recorded
\reviewmatch{Matches source input}
\noindent\textbf{Reviewer annotation}\par
\reviewerbox
\clearpage
\hypertarget{source-claim-8281b9d563b2f959}{}
\section*{Review row 23}
\textbf{Source locator:} {\footnotesize\raggedright cited publication, lines 1810-\allowbreak{}1835\par}
\paragraph{Verbatim source input}
\begin{ReviewVerbatim}
of a product is calculated to be
n
k
m
· +
·C
n+m n n+m
k + mC
=
,
n+m
which is equivalent to a prior-weighted rating system
with prior strength m and baseline prior distribution of
Beta(C, 1).
q=

(20)

From here, one may proceed with the proofs of the theorems in Section 3 so long as they can make statements about
the bias and variance terms in the above expression. In the
responsive model in this paper, N (qv , η) is a noisy random
variable that is a function of the play distribution in a multiarmed bandit. The theoretical challenge in proving stronger
results in the responsive setting is being able to characterize N (qv , η) for finite samples, and how it changes with η
– in this work, we show empirical results that the tradeoff
holds for a variety of consumer choice models (Thompson
sampling, and various models discussed in Appendix F).
\end{ReviewVerbatim}


\paragraph{Expanded PaperInterface specification}
\begin{ReviewVerbatim}
∀ (n k : ℕ) (m C : ℝ),
  MBJG25ProducerFairness.assumption_positive_rating_count n →
    have nR := ↑n;
    have kR := ↑k;
    nR / (nR + m) * (kR / nR) + m / (nR + m) * C = (kR + m * C) / (nR + m) ∧
      MBJG25ProducerFairness.paper_posterior_rating (m * C) (m * (1 - C)) k n = (kR + m * C) / (nR + m)
\end{ReviewVerbatim}

\paragraph{Lean proof endpoint (built separately)}\begin{ReviewVerbatim}
MBJG25ProducerFairness.paper_facing_eq20_bayesian_rating_equivalence
\end{ReviewVerbatim}

\paragraph{Recorded source-to-Spec screening}
\textbf{Verdict:} \texttt{not recorded}
\\
\textbf{Reason:} not recorded
\reviewmatch{Matches source input}
\noindent\textbf{Reviewer annotation}\par
\reviewerbox
\clearpage
\hypertarget{source-claim-48bd38a1dccb3b5f}{}
\section*{Review row 24}
\textbf{Source locator:} {\footnotesize\raggedright cited publication, lines 1903-\allowbreak{}2029\par}
\paragraph{Verbatim source input}
\begin{ReviewVerbatim}
Because of the ordinal nature of the ratings data in this
where the exact rating aggregation method is hidden to comsetting, we need to adjust our model primitives as defined in
bat fake reviews and brigading, sample mean rating systems
Section 2.1. Rather than a Beta-Bernoulli setting, we model
where the overall rating of a product is the average ratour ratings as being drawn from a categorical distribution
ing it has received over a certain period of time by trustwith a Dirichlet prior. Let K be the set of possible ratings
worthy reviewers, and Bayesian ratings, which calculates
for a product v (i.e. K = {1, 2, 3, 4, 5} for Amazon proda weighted average between the sample mean rating and a
ucts and K = {0, 1, 2, 3, 4, 5} for Goodreads products). If
population-level mean rating. Assuming the reviews are bi7
nary, both sample mean and Bayesian ratings are captured
Licensed under the MIT License
8
by our prior-weighted ratings system. Sample mean ratings
Licensed under the Apache 2.0 license


[form-feed in source extraction]
Product Category
Children
Comics & Graphic
Fantasy & Paranormal
History & Biography
Mystery, Crime, Thriller
Poetry
Romance
Young Adult
All Beauty
Beauty & Personal Care
Clothing, Shoes, Jewelry
Health & Personal Care
Industrial & Scientific
Magazine Subscriptions
Patio, Lawn, Garden
Sports & Outdoor
Tools, Home Improvement
Video Games

Platform
Goodreads
Goodreads
Goodreads
Goodreads
Goodreads
Goodreads
Goodreads
Goodreads
Amazon
Amazon
Amazon
Amazon
Amazon
Amazon
Amazon
Amazon
Amazon
Amazon

# Items
124.1k
89.4k
258.6k
302.9k
219.2k
36.5k
335.4k
93.4k
112.6k
94.3k
7.2mil
60.3k
427.5k
3.4k
851.7k
1.6mil
1.5mil
137.2k

# Ratings
734.6k
542.4k
3.4mil
2.1mil
1.8mil
154.6k
3.6mil
2.4mil
701.5k
2.1mil
66.0mil
494.1k
5.2mil
71.5k
16.5mil
19.6mil
27.0mil
4.6mil

Table 2: Table showing details about each dataset used in
Appendix E.

ratings for a product v are drawn from a categorical distribution with weights {pv,k }P
k∈K , then we define the true quality
qv of this product to be j∈K pv,j · j.
To estimate the quality of a product at time T , we define
S(v, k, T ) to be the set of k-star reviews for product v at
time T , and N (v, k, t) := |S(v, k, T )|. The prior parameters, instead of being for a Beta distribution, are now for a
Dirichlet distribution; they are given by α̂ = {α̂j }j∈K . Due
to the conjugacy of the Dirichlet distribution and the categorical distribution, we can define the estimated quality for
v at time t as

P

j∈K
q̂α̂ (v, t) = P

N (v, j, t) · j

j∈K N (v, j, t)

(21)

Our notions of prior strength and shape remain the same
as before; that is, if we fix a prior shape α̃ = {α̃j }j∈K ,
our choice of strength η ≥ 0 changes the Dirichlet prior
by setting α̂ = {η α̃j }j∈K . The rest of our also model definitions remain unchanged. We kept the same experimental
hyperparameters of the original experiment in Section 4.2,
with the exception of our η values now being chosen from
{0.001, 1, 10, 100, 1000, 10000}, and our Thompson sampling model for consumer choice now using Dirichlet priors rather than Beta priors. To fit the Dirichlet distribution to
\end{ReviewVerbatim}


\paragraph{Expanded PaperInterface specification}
\begin{ReviewVerbatim}
∀ {K : Type u_1} [inst : Fintype K] (ratingValue priorPseudoCount observedCount : K → ℝ),
  MBJG25ProducerFairness.paper_ordinal_posterior_rating ratingValue priorPseudoCount observedCount =
    (∑ j, (priorPseudoCount j + observedCount j) * ratingValue j) / ∑ j, (priorPseudoCount j + observedCount j)
\end{ReviewVerbatim}

\paragraph{Lean proof endpoint (built separately)}\begin{ReviewVerbatim}
MBJG25ProducerFairness.paper_ordinal_posterior_rating
\end{ReviewVerbatim}

\paragraph{Recorded source-to-Spec screening}
\textbf{Verdict:} \texttt{not recorded}
\\
\textbf{Reason:} not recorded
\reviewmatch{Matches source input}
\noindent\textbf{Reviewer annotation}\par
\reviewerbox
\clearpage
\hypertarget{source-claim-c141e929e91e7be9}{}
\section*{Review row 25}
\textbf{Source locator:} {\footnotesize\raggedright cited publication, lines 1903-\allowbreak{}2029\par}
\paragraph{Verbatim source input}
\begin{ReviewVerbatim}
Because of the ordinal nature of the ratings data in this
where the exact rating aggregation method is hidden to comsetting, we need to adjust our model primitives as defined in
bat fake reviews and brigading, sample mean rating systems
Section 2.1. Rather than a Beta-Bernoulli setting, we model
where the overall rating of a product is the average ratour ratings as being drawn from a categorical distribution
ing it has received over a certain period of time by trustwith a Dirichlet prior. Let K be the set of possible ratings
worthy reviewers, and Bayesian ratings, which calculates
for a product v (i.e. K = {1, 2, 3, 4, 5} for Amazon proda weighted average between the sample mean rating and a
ucts and K = {0, 1, 2, 3, 4, 5} for Goodreads products). If
population-level mean rating. Assuming the reviews are bi7
nary, both sample mean and Bayesian ratings are captured
Licensed under the MIT License
8
by our prior-weighted ratings system. Sample mean ratings
Licensed under the Apache 2.0 license


[form-feed in source extraction]
Product Category
Children
Comics & Graphic
Fantasy & Paranormal
History & Biography
Mystery, Crime, Thriller
Poetry
Romance
Young Adult
All Beauty
Beauty & Personal Care
Clothing, Shoes, Jewelry
Health & Personal Care
Industrial & Scientific
Magazine Subscriptions
Patio, Lawn, Garden
Sports & Outdoor
Tools, Home Improvement
Video Games

Platform
Goodreads
Goodreads
Goodreads
Goodreads
Goodreads
Goodreads
Goodreads
Goodreads
Amazon
Amazon
Amazon
Amazon
Amazon
Amazon
Amazon
Amazon
Amazon
Amazon

# Items
124.1k
89.4k
258.6k
302.9k
219.2k
36.5k
335.4k
93.4k
112.6k
94.3k
7.2mil
60.3k
427.5k
3.4k
851.7k
1.6mil
1.5mil
137.2k

# Ratings
734.6k
542.4k
3.4mil
2.1mil
1.8mil
154.6k
3.6mil
2.4mil
701.5k
2.1mil
66.0mil
494.1k
5.2mil
71.5k
16.5mil
19.6mil
27.0mil
4.6mil

Table 2: Table showing details about each dataset used in
Appendix E.

ratings for a product v are drawn from a categorical distribution with weights {pv,k }P
k∈K , then we define the true quality
qv of this product to be j∈K pv,j · j.
To estimate the quality of a product at time T , we define
S(v, k, T ) to be the set of k-star reviews for product v at
time T , and N (v, k, t) := |S(v, k, T )|. The prior parameters, instead of being for a Beta distribution, are now for a
Dirichlet distribution; they are given by α̂ = {α̂j }j∈K . Due
to the conjugacy of the Dirichlet distribution and the categorical distribution, we can define the estimated quality for
v at time t as

P

j∈K
q̂α̂ (v, t) = P

N (v, j, t) · j

j∈K N (v, j, t)

(21)

Our notions of prior strength and shape remain the same
as before; that is, if we fix a prior shape α̃ = {α̃j }j∈K ,
our choice of strength η ≥ 0 changes the Dirichlet prior
by setting α̂ = {η α̃j }j∈K . The rest of our also model definitions remain unchanged. We kept the same experimental
hyperparameters of the original experiment in Section 4.2,
with the exception of our η values now being chosen from
{0.001, 1, 10, 100, 1000, 10000}, and our Thompson sampling model for consumer choice now using Dirichlet priors rather than Beta priors. To fit the Dirichlet distribution to
\end{ReviewVerbatim}


\paragraph{Expanded PaperInterface specification}
\begin{ReviewVerbatim}
∀ {K : Type u_1} [inst : Fintype K] (ratingValue observedCount : K → ℝ),
  MBJG25ProducerFairness.paper_ordinal_posterior_rating ratingValue (fun x => 0) observedCount =
    (∑ j, observedCount j * ratingValue j) / ∑ j, observedCount j
\end{ReviewVerbatim}

\paragraph{Lean proof endpoint (built separately)}\begin{ReviewVerbatim}
MBJG25ProducerFairness.paper_facing_eq21_zero_prior_comparison
\end{ReviewVerbatim}

\paragraph{Recorded source-to-Spec screening}
\textbf{Verdict:} \texttt{not recorded}
\\
\textbf{Reason:} not recorded
\reviewmatch{Matches source input}
\noindent\textbf{Reviewer annotation}\par
\reviewerbox
\clearpage
\hypertarget{source-claim-c0449d6c0ae69374}{}
\section*{Review row 26}
\textbf{Source locator:} {\footnotesize\raggedright cited publication, lines 2090-\allowbreak{}2175\par}
\paragraph{Verbatim source input}
\begin{ReviewVerbatim}
F

Robustness tests for different selection
mechanisms

The responsive-setting empirical setup detailed in Section
4.2 assumes that agents are Thompson samplers, which
is meant to serve as a metaphor for real-world customers
roughly choosing the best product on the market based on
the ratings information available to them. The amount of
noise in this decision-making process can vary from platform to platform; for example, users in some online platforms may pay less attention to other user-generated reviews, due to the user interface de-emphasizing ratings. In
order to capture how our results may change in response to
this noisiness, we use this appendix to repeat our empirical
experiments while replacing our Thompson sampling choice
algorithm with lower-fidelity variants.
To generalize our Thompson sampling algorithm into
a parameterizable algorithm that we can tune, we use ksampling. This is a simple variant of Thompson sampling,
which is defined as follows: Instead of drawing a sample
from n different samples and selecting the arm that maximizes the expected reward conditional on that sample, we instead take the top-k most expected reward-maximizing arms
and select one of those arms uniformly at random. Note that
at k = 1, this approach is equivalent to traditional Thompson
sampling, while at k = n, we are picking an arm uniformly

1.0

Nucleus Size
1
2
3
4
5
Prior Strength
0.001
1
10
100
1000
10000

0.9
0.8
Expected Total Regret

4.0

0.7
0.6
0.5
0.4
0.3
0.2

0.100

0.125
0.150
0.175
0.200
Standard Deviation of Selection Rate

0.225

1
Figure 12: Scatterplot of average
standard deviation in selection rate against total expected regret for robustness tests
of the responsive setting when the consumer sampling algorithm is modified. Colours represent the nucleus size k,
while size represents the prior strength η. The efficiencyfairness trade-off illustrated in Figure 3 holds for all values
of k except for when k = 5, where sampling products purely
at random at each timestep incurs a regret-maximizing but
”fair” solution, and the value of η does not affect fairness or
efficiency.

at random.9
We repeat our KuaiRec experiments while replacing Thompson sampling with k-sampling, varying k ∈
{1, 2, 3, 4, 5}; we refer to k as the nucleus size. We choose
our η values from {0.001, 1, 10, 100, 1000, 10000}, and
keep all other hyperparameters unchanged from Section 4.2.
Note that because we keep the market size of 5 unchanged
from our original simulations, our sampling algorithm samples uniformly at random when k = 5.
We find that the central efficiency-fairness trade-off is surprisingly resilient to degradations in the consumer sampling
algorithm in identifying the best-performing product. While
the uniformly random selection case k = 5 behaves chaotically, due to the product choice at each timestep being completely randomly chosen and not informed by any kind of
prior knowledge, all other nucleus sizes yield an efficiencyversus-fairness curve that (while noisy for higher values of
k) qualitatively exhibit similar behavior as in our primary
experiments. In particular, as η increases, regret generally
increases while unfairness decreases.
\end{ReviewVerbatim}


\paragraph{Expanded PaperInterface specification}
\begin{ReviewVerbatim}
∀ {V : Type u_1} [inst : Fintype V] [DecidableEq V] [Nonempty V] (profile : V → ℝ) (k : ℕ),
  0 < k →
    k ≤ Fintype.card V →
      ∃ top,
        top.card = k ∧
          (∀ v ∈ top, ∀ w ∉ top, profile w ≤ profile v) ∧
            ∃ (htop : top.Nonempty), ∃ policy, policy = PMF.uniformOfFinset top htop
\end{ReviewVerbatim}

\paragraph{Lean proof endpoint (built separately)}\begin{ReviewVerbatim}
MBJG25ProducerFairness.paper_facing_k_sampling_mechanism
\end{ReviewVerbatim}

\paragraph{Recorded source-to-Spec screening}
\textbf{Verdict:} \texttt{not recorded}
\\
\textbf{Reason:} not recorded
\reviewmatch{Matches source input}
\noindent\textbf{Reviewer annotation}\par
\reviewerbox
\clearpage
\hypertarget{source-claim-03b433e77eab6bcb}{}
\section*{Review row 27}
\textbf{Source locator:} {\footnotesize\raggedright cited publication, lines 2090-\allowbreak{}2175\par}
\paragraph{Verbatim source input}
\begin{ReviewVerbatim}
F

Robustness tests for different selection
mechanisms

The responsive-setting empirical setup detailed in Section
4.2 assumes that agents are Thompson samplers, which
is meant to serve as a metaphor for real-world customers
roughly choosing the best product on the market based on
the ratings information available to them. The amount of
noise in this decision-making process can vary from platform to platform; for example, users in some online platforms may pay less attention to other user-generated reviews, due to the user interface de-emphasizing ratings. In
order to capture how our results may change in response to
this noisiness, we use this appendix to repeat our empirical
experiments while replacing our Thompson sampling choice
algorithm with lower-fidelity variants.
To generalize our Thompson sampling algorithm into
a parameterizable algorithm that we can tune, we use ksampling. This is a simple variant of Thompson sampling,
which is defined as follows: Instead of drawing a sample
from n different samples and selecting the arm that maximizes the expected reward conditional on that sample, we instead take the top-k most expected reward-maximizing arms
and select one of those arms uniformly at random. Note that
at k = 1, this approach is equivalent to traditional Thompson
sampling, while at k = n, we are picking an arm uniformly

1.0

Nucleus Size
1
2
3
4
5
Prior Strength
0.001
1
10
100
1000
10000

0.9
0.8
Expected Total Regret

4.0

0.7
0.6
0.5
0.4
0.3
0.2

0.100

0.125
0.150
0.175
0.200
Standard Deviation of Selection Rate

0.225

1
Figure 12: Scatterplot of average
standard deviation in selection rate against total expected regret for robustness tests
of the responsive setting when the consumer sampling algorithm is modified. Colours represent the nucleus size k,
while size represents the prior strength η. The efficiencyfairness trade-off illustrated in Figure 3 holds for all values
of k except for when k = 5, where sampling products purely
at random at each timestep incurs a regret-maximizing but
”fair” solution, and the value of η does not affect fairness or
efficiency.

at random.9
We repeat our KuaiRec experiments while replacing Thompson sampling with k-sampling, varying k ∈
{1, 2, 3, 4, 5}; we refer to k as the nucleus size. We choose
our η values from {0.001, 1, 10, 100, 1000, 10000}, and
keep all other hyperparameters unchanged from Section 4.2.
Note that because we keep the market size of 5 unchanged
from our original simulations, our sampling algorithm samples uniformly at random when k = 5.
We find that the central efficiency-fairness trade-off is surprisingly resilient to degradations in the consumer sampling
algorithm in identifying the best-performing product. While
the uniformly random selection case k = 5 behaves chaotically, due to the product choice at each timestep being completely randomly chosen and not informed by any kind of
prior knowledge, all other nucleus sizes yield an efficiencyversus-fairness curve that (while noisy for higher values of
k) qualitatively exhibit similar behavior as in our primary
experiments. In particular, as η increases, regret generally
increases while unfairness decreases.
\end{ReviewVerbatim}


\paragraph{Expanded PaperInterface specification}
\begin{ReviewVerbatim}
∀ {V : Type u_1} [Fintype V] [DecidableEq V] [Nonempty V] (profile : V → ℝ),
  ∃ top,
    top.card = 1 ∧
      ∃ v,
        top = {v} ∧
          (∀ (w : V), profile w ≤ profile v) ∧ ∃ (htop : top.Nonempty), PMF.uniformOfFinset top htop = PMF.pure v
\end{ReviewVerbatim}

\paragraph{Lean proof endpoint (built separately)}\begin{ReviewVerbatim}
MBJG25ProducerFairness.paper_facing_k_sampling_one_is_argmax
\end{ReviewVerbatim}

\paragraph{Recorded source-to-Spec screening}
\textbf{Verdict:} \texttt{not recorded}
\\
\textbf{Reason:} not recorded
\reviewmatch{Matches source input}
\noindent\textbf{Reviewer annotation}\par
\reviewerbox
\clearpage
\hypertarget{source-claim-48e4400514808bc7}{}
\section*{Review row 28}
\textbf{Source locator:} {\footnotesize\raggedright cited publication, lines 2090-\allowbreak{}2175\par}
\paragraph{Verbatim source input}
\begin{ReviewVerbatim}
F

Robustness tests for different selection
mechanisms

The responsive-setting empirical setup detailed in Section
4.2 assumes that agents are Thompson samplers, which
is meant to serve as a metaphor for real-world customers
roughly choosing the best product on the market based on
the ratings information available to them. The amount of
noise in this decision-making process can vary from platform to platform; for example, users in some online platforms may pay less attention to other user-generated reviews, due to the user interface de-emphasizing ratings. In
order to capture how our results may change in response to
this noisiness, we use this appendix to repeat our empirical
experiments while replacing our Thompson sampling choice
algorithm with lower-fidelity variants.
To generalize our Thompson sampling algorithm into
a parameterizable algorithm that we can tune, we use ksampling. This is a simple variant of Thompson sampling,
which is defined as follows: Instead of drawing a sample
from n different samples and selecting the arm that maximizes the expected reward conditional on that sample, we instead take the top-k most expected reward-maximizing arms
and select one of those arms uniformly at random. Note that
at k = 1, this approach is equivalent to traditional Thompson
sampling, while at k = n, we are picking an arm uniformly

1.0

Nucleus Size
1
2
3
4
5
Prior Strength
0.001
1
10
100
1000
10000

0.9
0.8
Expected Total Regret

4.0

0.7
0.6
0.5
0.4
0.3
0.2

0.100

0.125
0.150
0.175
0.200
Standard Deviation of Selection Rate

0.225

1
Figure 12: Scatterplot of average
standard deviation in selection rate against total expected regret for robustness tests
of the responsive setting when the consumer sampling algorithm is modified. Colours represent the nucleus size k,
while size represents the prior strength η. The efficiencyfairness trade-off illustrated in Figure 3 holds for all values
of k except for when k = 5, where sampling products purely
at random at each timestep incurs a regret-maximizing but
”fair” solution, and the value of η does not affect fairness or
efficiency.

at random.9
We repeat our KuaiRec experiments while replacing Thompson sampling with k-sampling, varying k ∈
{1, 2, 3, 4, 5}; we refer to k as the nucleus size. We choose
our η values from {0.001, 1, 10, 100, 1000, 10000}, and
keep all other hyperparameters unchanged from Section 4.2.
Note that because we keep the market size of 5 unchanged
from our original simulations, our sampling algorithm samples uniformly at random when k = 5.
We find that the central efficiency-fairness trade-off is surprisingly resilient to degradations in the consumer sampling
algorithm in identifying the best-performing product. While
the uniformly random selection case k = 5 behaves chaotically, due to the product choice at each timestep being completely randomly chosen and not informed by any kind of
prior knowledge, all other nucleus sizes yield an efficiencyversus-fairness curve that (while noisy for higher values of
k) qualitatively exhibit similar behavior as in our primary
experiments. In particular, as η increases, regret generally
increases while unfairness decreases.
\end{ReviewVerbatim}


\paragraph{Expanded PaperInterface specification}
\begin{ReviewVerbatim}
∀ {V : Type u_1} [inst : Fintype V] [DecidableEq V] [inst_2 : Nonempty V] (profile : V → ℝ),
  ∃ top, top.card = Fintype.card V ∧ ∃ (htop : top.Nonempty), PMF.uniformOfFinset top htop = PMF.uniformOfFintype V
\end{ReviewVerbatim}

\paragraph{Lean proof endpoint (built separately)}\begin{ReviewVerbatim}
MBJG25ProducerFairness.paper_facing_k_sampling_market_size_is_uniform
\end{ReviewVerbatim}

\paragraph{Recorded source-to-Spec screening}
\textbf{Verdict:} \texttt{not recorded}
\\
\textbf{Reason:} not recorded
\reviewmatch{Matches source input}
\noindent\textbf{Reviewer annotation}\par
\reviewerbox

\end{document}
